\documentclass[a4paper,oneside,11pt]{book}

% Package yang diperlukan
\usepackage{graphicx}
\usepackage{titling}
\usepackage{geometry}
\usepackage{listings}
\usepackage{xcolor}
\usepackage{float}
\usepackage{indentfirst}
\usepackage{array}
\usepackage[table]{xcolor}
\usepackage{tabularx}
\usepackage[indonesian]{babel}
\usepackage{multirow}
\usepackage{booktabs}
\usepackage{url}
\usepackage{mathptmx}   
\usepackage{helvet}    
\usepackage{amsfonts}   
\usepackage{tikz}
\usepackage{amsmath}
\usepackage{xurl}
\usepackage[breaklinks=true]{hyperref}
\usetikzlibrary{shapes.geometric, arrows}

% Untuk manajemen font dan karakter khusus
\usepackage[T1]{fontenc}       % Encoding font Latin modern
\usepackage{textcomp}          % Simbol tambahan (text companion)

% Untuk referensi silang yang lebih baik
\usepackage[all]{hypcap}       % Perbaiki link ke gambar/tabel

% Untuk kontrol paragraf
\usepackage{parskip}           % Spasi antar paragraf konsisten

% Untuk pengaturan halaman
\usepackage{fancyhdr}          % Header/footer custom
\pagestyle{fancy}              % Aktifkan style fancy
\fancyhf{}                     % Kosongkan header/footer
\rfoot{\thepage}               % Nomor halaman di kanan bawah

% Untuk gambar yang lebih fleksibel
\usepackage{adjustbox}         % Penyesuaian ukuran gambar otomatis

% Untuk kode program tambahan
\usepackage{etoolbox}          % Tools untuk programming LaTeX

\hypersetup{
    colorlinks=true,
    linkcolor=black,
    citecolor=blue,
    urlcolor=blue
}

% Pengaturan margin
\geometry{
    left=3cm,
    right=2.5cm,
    top=3cm,
    bottom=3cm
}

% Pengaturan untuk kode HTML/Python
\definecolor{codegreen}{rgb}{0,0.6,0}
\definecolor{codegray}{rgb}{0.5,0.5,0.5}
\definecolor{codepurple}{rgb}{0.58,0,0.82}
\definecolor{backcolour}{rgb}{0.95,0.95,0.92}

\lstdefinestyle{htmlstyle}{
    backgroundcolor=\color{backcolour},   
    commentstyle=\color{codegreen},
    keywordstyle=\color{codepurple},
    numberstyle=\tiny\color{codegray},
    stringstyle=\color{codegreen},
    basicstyle=\ttfamily\footnotesize,
    breakatwhitespace=false,         
    breaklines=true,                 
    captionpos=b,                    
    keepspaces=true,                 
    numbers=left,                    
    numbersep=5pt,                  
    showspaces=false,                
    showstringspaces=false,
    showtabs=false,                  
    tabsize=2,
    frame=single
}

\begin{document}

% ================= HALAMAN JUDUL =================
\begin{titlingpage}
\begin{center}
% Spasi atas
\vspace*{0.5cm}

% Bagian Judul Atas
{\large \textbf{LAPORAN PRAKTIKUM}} \\[0.3cm]
{\large \textbf{PRAKTIKUM PENAMBANGAN DATA}} \\[0.3cm]
{\large \textbf{PERTEMUAN 5}} \\[0.3cm]
{\large \textbf{FORECASTING}} \\[2.5cm] % Ubah judul materi sesuai kebutuhan

% Bagian Logo
% Pastikan file lambangugm.png berada di satu folder dengan file .tex
\includegraphics[width=6cm]{lambangugm.png} \\[2.5cm]

% Bagian Identitas
\normalsize
Disusun oleh:\\[0.5cm]
\begin{tabular}{@{}l@{\hspace{1cm}}l@{}}
Nama           & : Hilarius Christiano Avin Paliling \\
NIM            & : 24/542159/SV/25009 \\
Kelas          & : PL5B1 \\
Dosen Pengampu & : Dr.Eng. Ir. Ganjar Alfian, S.T., M.Eng.
\end{tabular}

% vfill akan secara otomatis mendorong teks di bawahnya ke dasar halaman
\vfill 

% Bagian Institusi Bawah
{\large \textbf{PROGRAM STUDI D-IV TEKNOLOGI REKAYASA PERANGKAT LUNAK}} \\[0.1cm]
{\large \textbf{DEPARTEMEN TEKNIK ELEKTRO DAN INFORMATIKA}} \\[0.1cm]
{\large \textbf{SEKOLAH VOKASI}} \\[0.1cm]
{\large \textbf{UNIVERSITAS GADJAH MADA}} \\[0.1cm]
{\large \textbf{YOGYAKARTA}} \\[0.1cm]
{\large \textbf{2026}}

% Spasi batas bawah agar tidak menyentuh margin
\vspace*{0.5cm}
\end{center}
\end{titlingpage}
% =================================================

% Daftar Isi
\tableofcontents
\listoffigures
\newpage

\chapter{Tujuan Praktikum}
\begin{enumerate}
    \item Mahasiswa dapat menerapkan metode forecasting berbasis machine learning, meliputi Neural Network, K-Nearest Neighbors (KNN), Decision Tree, Support Vector Regression (SVR), dan Random Forest, untuk memprediksi nilai periode berikutnya berdasarkan data historis
    \item Mahasiswa dapat melakukan analisis dan perbandingan performa model menggunakan metrik evaluasi Root Mean Squared Error (RMSE) dan Pearson Correlation Coefficient
    \item Mahasiswa dapat menentukan model forecasting yang paling akurat berdasarkan hasil evaluasi menggunakan RMSE dan Pearson Correlation Coefficient
\end{enumerate}

\chapter{Dasar Teori}

\section{Supervised Learning}
\textit{Supervised learning} merupakan salah satu paradigma fundamental dalam \textit{machine learning} di mana algoritma dilatih menggunakan kumpulan data berlabel (\textit{labeled data}) agar model mampu mengenali karakteristik dan pola hubungan antara variabel masukan (\textit{input features}) dengan variabel target luaran guna melakukan estimasi prediksi atau klasifikasi [\ref{retnoningsih}]. Pada paradigma \textit{supervised learning}, sistem diberikan masukan beserta luaran target yang diinginkan; sistem kemudian mempelajari pola data historis tersebut sebagai acuan inferensi pada data baru di masa mendatang [\ref{santoso}].

\section{Forecasting (Peramalan Deret Waktu)}
\textit{Forecasting} adalah proses memproyeksikan atau memprediksi suatu keadaan di masa kini dan masa depan dengan menguji dan mengevaluasi pola keadaan di masa lalu [\ref{vimala}]. Peramalan memegang peranan krusial bagi organisasi maupun industri bisnis dalam mendukung proses pengambilan keputusan manajemen yang tepat [\ref{maysofa}]. Berdasarkan karakteristik keluarannya, peramalan dapat dibedakan menjadi:
\begin{enumerate}
    \item \textbf{Ramalan Kualitatif:} Bersifat deskriptif dan tidak berbentuk angka, seperti perkiraan cuaca cerah atau sentimen pasar secara umum.
    \item \textbf{Ramalan Kuantitatif:} Dinyatakan dalam bentuk angka atau besaran numerik kontinu berdasarkan formulasi matematis terhadap data historis.
\end{enumerate}
Dalam konteks data deret waktu (\textit{time series}), peramalan memanfaatkan ketergantungan temporal antarpengamatan berurutan untuk memperkirakan nilai target periode ke depan ($t+h$) berdasarkan riwayat pengamatan masa lalu ($t-1, t-2, \dots, t-n$) [\ref{modulPraktikum}].

\section{Metode Sliding Window dan Rekayasa Fitur Statistik}
Untuk menerapkan algoritma \textit{supervised learning} pada data sekuensial deret waktu, diperlukan tahapan rekayasa fitur khusus sesuai modul praktikum:
\begin{enumerate}
    \item \textbf{Teknik Jendela Geser (\textit{Sliding Window}):}\\
    Metode \textit{sliding window} mentransformasikan urutan sekuensial satu dimensi menjadi pasangan matriks fitur masukan ($X$) dan vektor target ($y$) [\ref{modulPraktikum}]. Jendela berukuran tetap $n\_steps\_in$ digeser sepanjang deret waktu dengan \textit{stride} satu satuan waktu untuk mengumpulkan nilai historis sebagai variabel bebas, sementara pengamatan pada langkah ke-($n\_steps\_in + n\_steps\_out$) ditetapkan sebagai target:
    \begin{equation}
        X_i = [s_i, s_{i+1}, \dots, s_{i+n\_steps\_in-1}], \quad y_i = s_{i+n\_steps\_in+n\_steps\_out-1}
    \end{equation}
    Pemisahan data latih dan data uji wajib mempertahankan kronologi waktu secara urut (\texttt{shuffle=False}) guna mencegah terjadinya kebocoran informasi masa depan (\textit{data leakage}).

    \item \textbf{Ekstraksi Fitur Statistik Deskriptif:}\\
    Guna memperkaya representasi distribusi data dalam jendela waktu tertentu dan meningkatkan performa model tanpa memerlukan jendela \textit{lag} yang terlalu panjang, diekstraksi 8 fitur statistik deskriptif dari setiap sampel jendela masukan [\ref{modulPraktikum}]:
    \begin{itemize}
        \item Nilai minimum (\texttt{min}) dan nilai maksimum (\texttt{max}).
        \item Selisih rentang (\texttt{diff} = $\max - \min$).
        \item Standar deviasi (\texttt{std}) dan rata-rata aritmatika (\texttt{mean}).
        \item Nilai tengah (\texttt{median}).
        \item Tingkat keruncingan distribusi (\texttt{kurtosis}) dan tingkat kemiringan distribusi (\texttt{skewness}).
    \end{itemize}
\end{enumerate}

\section{Algoritma Machine Learning untuk Forecasting}
Sesuai dengan panduan modul praktikum, empat algoritma \textit{supervised learning} digunakan untuk membangun model peramalan deret waktu:
\begin{enumerate}
    \item \textbf{Multilayer Perceptron (MLP):} Model jaringan saraf tiruan umpan maju (\textit{feedforward artificial neural network}) yang terdiri dari lapisan masukan, satu atau lebih lapisan tersembunyi (\textit{hidden layers}), dan lapisan keluaran. MLP mampu memodelkan aproksimasi fungsi non-linear kompleks yang menghubungkan riwayat \textit{lag} dan fitur statistik dengan target peramalan melalui mekanisme propagasi balik (\textit{backpropagation}).
    \item \textbf{K-Nearest Neighbors (KNN):} Algoritma berbasis memori non-parametrik (\textit{instance-based learning}) yang mengestimasi nilai target suatu sampel berdasarkan rata-rata nilai dari $k$ tetangga terdekatnya dalam ruang fitur menggunakan perhitungan jarak geometris (\textit{Euclidean distance}).
    \item \textbf{Decision Tree (DT):} Model pohon keputusan regresi yang mempartisi ruang fitur secara rekursif menjadi sub-wilayah yang lebih homogen berdasarkan kriteria minimalisasi variansi galat (\textit{Mean Squared Error}).
    \item \textbf{Random Forest (RF):} Algoritma ansambel (\textit{ensemble learning}) berbasis \textit{bagging} (\textit{bootstrap aggregating}) yang menggabungkan hasil estimasi dari kumpulan pohon keputusan independen. Mekanisme ini mereduksi variansi model secara signifikan, menjadikannya sangat tangguh terhadap derau (\textit{noise}) data deret waktu dan meminimalkan risiko \textit{overfitting}.
\end{enumerate}

\clearpage
\section{Metrik Evaluasi Model Forecasting}
Evaluasi model peramalan dilakukan menggunakan dua metrik komplementer untuk mengukur besarnya galat numerik sekaligus keselarasan tren prediksi terhadap nilai aktual [\ref{modulPraktikum}]:
\begin{enumerate}
    \item \textbf{Root Mean Squared Error (RMSE):}\\
    RMSE mengukur akar kuadrat dari rata-rata selisih kuadrat antara nilai aktual ($y_{i}$) dan nilai hasil prediksi ($\hat{y}_{i}$) [\ref{chai}]:
    \begin{equation}
        RMSE = \sqrt{\frac{1}{n}\sum_{i=1}^{n}\left(y_{i} - \hat{y}_{i}\right)^2}
    \end{equation}
    RMSE memberikan penalti yang besar pada kesalahan ekstrem sehingga sangat sensitif terhadap galat besar. Nilai RMSE yang lebih kecil menandakan performa model yang semakin baik karena memiliki tingkat kesalahan prediksi yang rendah.

    \item \textbf{Pearson Correlation Coefficient ($r$):}\\
    Koefisien korelasi Pearson mengukur derajat keeratan hubungan linear dan keselarasan pola pergerakan antara nilai aktual ($y$) dan nilai prediksi ($\hat{y}$) [\ref{benesty}]:
    \begin{equation}
        r = \frac{\sum_{i=1}^{n}(y_{i} - \bar{y})(\hat{y}_{i} - \bar{\hat{y}})}{\sqrt{\sum_{i=1}^{n}(y_{i} - \bar{y})^2 \sum_{i=1}^{n}(\hat{y}_{i} - \bar{\hat{y}})^2}}
    \end{equation}
    Koefisien ini memiliki rentang nilai antara $-1$ hingga $1$. Nilai yang mendekati $1$ menunjukkan bahwa kurva prediksi model sangat baik dalam menangkap dan mengikuti dinamika tren pola data aktual.
\end{enumerate}

\chapter{Hasil dan Pembahasan}

\section{Hasil Praktikum}
Pada bagian Hasil Praktikum, dilakukan eksperimen peramalan (\textit{forecasting}) data deret waktu (\textit{time series}) menggunakan \textit{Bike Sharing Dataset}. Eksperimen ini bertujuan memprediksi jumlah peminjaman sepeda (\texttt{cnt}) untuk satu periode ke depan ($t+1$) berdasarkan data historis 7 hari sebelumnya ($t-7$ hingga $t-1$) dengan memanfaatkan metode \textit{sliding window}, penambahan fitur statistik deret waktu, serta pemodelan menggunakan empat algoritma \textit{supervised learning} (\textit{Multilayer Perceptron}, \textit{K-Nearest Neighbors}, \textit{Decision Tree}, dan \textit{Random Forest}).

\subsection{Import Library}
Tahap awal dimulai dengan mempersiapkan seluruh modul dan pustaka Python yang dibutuhkan untuk manipulasi data, rekayasa fitur statistik, pembuatan model \textit{machine learning}, dan evaluasi performa peramalan.

\begin{figure}[H]
    \centering
    \includegraphics[width=0.55\textwidth]{1.png}
    \caption{Source Code Import Library}
    \label{fig:1}
\end{figure}
Penjelasan Gambar :
\begin{itemize}
    \item \textbf{Input:} Mengimpor pustaka dasar manipulasi data numerik dan tabel (\texttt{numpy}, \texttt{pandas}), pustaka visualisasi grafik (\texttt{matplotlib.pyplot}), algoritma regresi dari \texttt{sklearn} (\texttt{MLPRegressor}, \texttt{KNeighborsRegressor}, \texttt{DecisionTreeRegressor}, \texttt{RandomForestRegressor}, \texttt{SVR}), fungsi evaluasi metrik kesalahan (\texttt{mean\_squared\_error}, \texttt{sqrt}), fungsi pemisahan data (\texttt{train\_test\_split}), serta fungsi statistik dari pustaka \texttt{scipy.stats} (\texttt{pearsonr}, \texttt{kurtosis}, \texttt{skew}).
    \item \textbf{Output:} Seluruh pustaka berhasil dimuat ke dalam memori lingkungan kerja Google Colab tanpa terjadi kesalahan dependensi.
\end{itemize}

\subsection{Pemuatan Dataset dan Penyesuaian Format Data}
Dataset \textit{Bike Sharing} dimuat secara langsung melalui tautan repositori daring, kemudian dilakukan pemisahan atribut target dan penyesuaian tipe data.

\begin{figure}[H]
    \centering
    \includegraphics[width=0.40\textwidth]{2a.png}
    \caption{Source Code dan Output Load Dataset Bike Sharing}
    \label{fig:2a}
    
    \vspace{0.1cm}
    
    \includegraphics[width=0.3\textwidth]{2b.png}
    \caption{Source Code dan Output Pemilihan Fitur Target}
    \label{fig:2b}
\end{figure}
Penjelasan Gambar :
\begin{itemize}
    \item \textbf{Input:} Pada Gambar \ref{fig:2a}, dataset \texttt{bikesharing\_day.csv} dibaca menggunakan \texttt{pd.read\_csv()}, kolom tanggal \texttt{dteday} dikonversi ke tipe \texttt{datetime}, dan fungsi \texttt{df.head()} dipanggil untuk memeriksa struktur data awal. Pada Gambar \ref{fig:2b}, atribut target peminjaman sepeda (\texttt{cnt}) diekstrak dan diduplikasi ke dalam bentuk matriks 2 dimensi \texttt{df\_X = df\_ori[['cnt', 'cnt']]}, kemudian dikonversi menjadi larik bertipe \textit{floating-point} via \texttt{astype(float).values}.
    \item \textbf{Output:} Gambar \ref{fig:2a} menampilkan 5 sampel teratas dataset yang memuat 16 atribut (termasuk kondisi cuaca, temperatur, kelembapan, dan total peminjaman). Gambar \ref{fig:2b} menampilkan sampel larik data \texttt{cnt} yang siap diproses ke tahap pembentukan jendela waktu.
\end{itemize}

\subsection{Definisi Fungsi Sliding Window}
Untuk mengonversi data deret waktu sekuensial satu dimensi menjadi format pasangan fitur input ($X$) dan target ($y$) yang kompatibel dengan algoritma \textit{supervised learning}, dibuat sebuah fungsi modular \textit{sliding window}.

\begin{figure}[H]
    \centering
    \includegraphics[width=0.55\textwidth]{3.png}
    \caption{Source Code Definisi Fungsi Sliding Window}
    \label{fig:3}
\end{figure}
Penjelasan Gambar :
\begin{itemize}
    \item \textbf{Input:} Mendefinisikan fungsi \texttt{split\_sequences(sequences, n\_steps\_in, n\_steps\_out)} yang menerima masukan berupa deret data (\texttt{sequences}), ukuran jendela input historis (\texttt{n\_steps\_in}), dan langkah periode target ke depan (\texttt{n\_steps\_out}). Fungsi melakukan iterasi secara dinamis untuk memotong segmen data masa lalu (\texttt{seq\_x}) sebagai prediktor dan nilai pada horizon ke depan (\texttt{seq\_y}) sebagai label keluaran.
    \item \textbf{Output:} Fungsi mengembalikan pasangan larik NumPy \texttt{array(X)} berdimensi tiga dan \texttt{array(y)} berdimensi satu yang merepresentasikan struktur data deret waktu berbasis jendela geser.
\end{itemize}

\subsection{Penambahan Fitur Statistik}
Guna memperkaya representasi karakteristik distribusi nilai dalam setiap jendela geser dan meningkatkan akurasi prediksi model, disusun fungsi pengekstraksi fitur statistik ringkasan.

\begin{figure}[H]
    \centering
    \includegraphics[width=0.35\textwidth]{4.png}
    \caption{Source Code Fungsi Penambahan Fitur Statistik}
    \label{fig:4}
\end{figure}
Penjelasan Gambar :
\begin{itemize}
    \item \textbf{Input:} Mendefinisikan fungsi \texttt{stats\_features(input\_data)} yang melakukan iterasi pada setiap jendela sampel data input untuk menghitung 8 indikator statistik deskriptif: nilai minimum (\texttt{np.min}), nilai maksimum (\texttt{np.max}), rentang selisih (\texttt{diff = max - min}), standar deviasi (\texttt{np.std}), rata-rata aritmatika (\texttt{np.mean}), nilai tengah (\texttt{np.median}), tingkat keruncingan kurva (\texttt{kurtosis}), dan tingkat kemiringan distribusi (\texttt{skew}). Kedelapan fitur ini kemudian digabungkan (\texttt{np.append}) ke dalam vektor fitur asli.
    \item \textbf{Output:} Fungsi menghasilkan matriks fitur baru yang telah diperluas dimensinya (\textit{augmented features}) dengan bentuk akhir \texttt{inp.reshape(len(input\_data), -1)}.
\end{itemize}

\subsection{Prapemrosesan Data (Data Preprocessing)}
Tahap prapemrosesan mencakup pembentukan jendela waktu dengan parameter $n\_steps\_in = 7$ dan $n\_steps\_out = 1$, perataan dimensi (\textit{flattening}), pembagian data latih dan uji kronologis, penambahan fitur statistik, serta inspeksi visual distribusi deret waktu.

\begin{figure}[H]
    \centering
    \includegraphics[width=0.55\textwidth]{5a.png}
    \caption{Source Code dan Eksekusi Pipeline Prapemrosesan Data}
    \label{fig:5a}
    
    \vspace{0.1cm}
    
    \includegraphics[width=0.35\textwidth]{5b.png}
    \caption{Output Verifikasi Matriks Fitur dan Target}
    \label{fig:5b}
\end{figure}

\begin{figure}[H]
    \centering
    \includegraphics[width=0.35\textwidth]{5c.png}
    \caption{Output Struktur Data Time Series Terindeks}
    \label{fig:5c}
    
    \vspace{0.3cm}
    
    \includegraphics[width=0.55\textwidth]{5d.png}
    \caption{Plot Visualisasi Tren Data Historis Peminjaman Sepeda}
    \label{fig:5d}
\end{figure}
Penjelasan Gambar :
\begin{itemize}
    \item \textbf{Input:} Pada Gambar \ref{fig:5a}, parameter \texttt{n\_steps\_in = 7} dan \texttt{n\_steps\_out = 1} diterapkan pada fungsi \texttt{split\_sequences}. Matriks $X$ kemudian diratakan menjadi 2 dimensi dan dipisahkan menjadi data latih (80\%) dan data uji (20\%) menggunakan fungsi \texttt{train\_test\_split} dengan parameter \texttt{shuffle=False} guna mempertahankan integritas urutan waktu temporal. Selanjutnya, fungsi \texttt{stats\_features} diterapkan secara terpisah pada \texttt{X\_train} dan \texttt{X\_test}. Pada Gambar \ref{fig:5d}, dibuat plot visualisasi pergerakan total peminjaman sepeda (\texttt{cnt}) terhadap indeks waktu menggunakan \texttt{plt.plot()}.
    \item \textbf{Output:} Gambar \ref{fig:5b} memvalidasi dimensi larik data hasil transformasi: $X$ memiliki 724 baris sampel dengan 7 fitur lag awal, sampel pertama $X[0]$ memuat nilai peminjaman 7 hari pertama, dan $y[0] = 959.0$ merupakan target hari ke-8. Gambar \ref{fig:5c} menampilkan pratinjau tabel 731 baris terindeks tanggal. Gambar \ref{fig:5d} memperlihatkan pola musiman (\textit{seasonality}) dan tren kenaikan volume peminjaman sepeda dari awal tahun 2011 hingga akhir tahun 2012.
\end{itemize}

\subsection{Pelatihan Model Machine Learning}
Empat algoritma \textit{machine learning} dilatih menggunakan arsitektur fungsi terisolasi yang menerima parameter data latih dan data uji, kemudian melakukan prediksi sekaligus menghitung metrik evaluasi secara otomatis.

\begin{figure}[H]
    \centering
    \includegraphics[width=0.55\textwidth]{6ab.png}
    \caption{Source Code Fungsi Pelatihan Model MLP dan KNN}
    \label{fig:6ab}
    
    \vspace{0.1cm}
    
    \includegraphics[width=0.55\textwidth]{6cd.png}
    \caption{Source Code Fungsi Pelatihan Model Decision Tree dan Random Forest}
    \label{fig:6cd}
\end{figure}
Penjelasan Gambar :
\begin{itemize}
    \item \textbf{Input:} Pada Gambar \ref{fig:6ab}, fungsi \texttt{mlp()} menginisialisasi \texttt{MLPRegressor(random\_state=42)} berbasis jaringan saraf tiruan, sedangkan fungsi \texttt{knn()} menginisialisasi \texttt{KNeighborsRegressor()}. Pada Gambar \ref{fig:6cd}, fungsi \texttt{dt()} menginisialisasi pohon keputusan \texttt{DecisionTreeRegressor(random\_state=42)}, dan fungsi \texttt{rf()} menginisialisasi algoritma ensemble \texttt{RandomForestRegressor(random\_state=42)}. Setiap fungsi melakukan fitting model (\texttt{model.fit(X\_train, y\_train)}), menghasilkan prediksi \texttt{y\_pred}, serta menghitung nilai RMSE dan koefisien korelasi Pearson.
    \item \textbf{Output:} Seluruh fungsi pemodelan terkompilasi dan siap dipanggil untuk proses pelatihan serta inferensi serentak.
\end{itemize}

\subsection{Pengujian dan Prediksi Data Uji}
Seluruh fungsi model yang telah didefinisikan dieksekusi secara berurutan menggunakan data uji \texttt{X\_test} dan \texttt{y\_test}.

\begin{figure}[H]
    \centering
    \includegraphics[width=0.65\textwidth]{7.png}
    \caption{Source Code Eksekusi Pelatihan dan Prediksi Model}
    \label{fig:7}
\end{figure}
Penjelasan Gambar :
\begin{itemize}
    \item \textbf{Input:} Menjalankan pemanggilan fungsi pemodelan \texttt{mlp()}, \texttt{knn()}, \texttt{dt()}, dan \texttt{rf()} dengan melewatkan argumen \texttt{X\_train}, \texttt{X\_test}, \texttt{y\_train}, dan \texttt{y\_test}.
    \item \textbf{Output:} Masing-masing fungsi berhasil mengembalikan nilai metrik RMSE, nilai korelasi Pearson, serta larik hasil prediksi (\texttt{y\_pred\_mlp}, \texttt{y\_pred\_knn}, \texttt{y\_pred\_dt}, dan \texttt{y\_pred\_rf}) untuk seluruh sampel data uji.
\end{itemize}

\subsection{Visualisasi Hasil Prediksi Model}
Untuk mengevaluasi secara kualitatif seberapa presisi pola peramalan tiap model dalam mengikuti fluktuasi data riil, dilakukan visualisasi perbandingan kurva prediksi terhadap nilai aktual (\texttt{y\_test}).

\begin{figure}[H]
    \centering
    \includegraphics[width=0.55\textwidth]{8a.png}
    \caption{Plot Visualisasi Hasil Prediksi Model Multilayer Perceptron (MLP)}
    \label{fig:8a}
    
    \vspace{0.3cm}
    
    \includegraphics[width=0.55\textwidth]{8b.png}
    \caption{Plot Visualisasi Hasil Prediksi Model K-Nearest Neighbors (KNN)}
    \label{fig:8b}
\end{figure}

\begin{figure}[H]
    \centering
    \includegraphics[width=0.55\textwidth]{8c.png}
    \caption{Plot Visualisasi Hasil Prediksi Model Decision Tree (DT)}
    \label{fig:8c}
    
    \vspace{0.3cm}
    
    \includegraphics[width=0.55\textwidth]{8d.png}
    \caption{Plot Visualisasi Hasil Prediksi Model Random Forest (RF)}
    \label{fig:8d}
\end{figure}
Penjelasan Gambar :
\begin{itemize}
    \item \textbf{Input:} Perintah pembuatan grafik deret waktu dengan pustaka \texttt{matplotlib.pyplot} yang memplotkan kurva data riil (\texttt{Real data}) bersanding dengan kurva hasil estimasi masing-masing algoritma pada rentang sumbu waktu $t$.
    \item \textbf{Output:} Gambar \ref{fig:8a} memperlihatkan kurva MLP yang sangat adaptif dan halus dalam menangkap pola kenaikan maupun penurunan volume peminjaman sepeda. Gambar \ref{fig:8b} (KNN) memperlihatkan prediksi yang cukup baik namun memiliki kecenderungan efek tangga (\textit{step-like function}). Gambar \ref{fig:8c} (Decision Tree) menunjukkan fluktuasi kaku yang rentan terhadap galat ekstrem pada lembah data. Gambar \ref{fig:8d} (Random Forest) menunjukkan kurva yang jauh lebih stabil dan akurat berkat mekanisme perataan dari kumpulan pohon keputusan.
\end{itemize}

\clearpage
\subsection{Evaluasi Performa Model}
Kinerja kuantitatif keempat algoritma diuji berdasarkan metrik kesalahan rata-rata kuadrat (RMSE) dan tingkat keeratan hubungan linear kurva prediksi dengan data aktual (\textit{Pearson Correlation Coefficient}).

\begin{figure}[H]
    \centering
    \includegraphics[width=0.6\textwidth]{9.png}
    \caption{Output Evaluasi Metrik Performa Model Forecasting}
    \label{fig:9}
\end{figure}
Penjelasan Gambar :
\begin{itemize}
    \item \textbf{Input:} Perintah pencetakan nilai evaluasi RMSE dan korelasi Pearson untuk model MLP, KNN, Decision Tree, dan Random Forest menggunakan format string terstruktur.
    \item \textbf{Output:} Menampilkan log performa numerik sebagai berikut:
    \begin{itemize}
        \item \textbf{MLP:} RMSE = $1190.549$, Korelasi Pearson = $0.778$
        \item \textbf{KNN:} RMSE = $1357.710$, Korelasi Pearson = $0.696$
        \item \textbf{Decision Tree (DT):} RMSE = $1854.465$, Korelasi Pearson = $0.484$
        \item \textbf{Random Forest (RF):} RMSE = $1324.452$, Korelasi Pearson = $0.716$
    \end{itemize}
    Berdasarkan hasil kuantitatif tersebut, model \textbf{Multilayer Perceptron (MLP)} mencatatkan performa terbaik dengan nilai RMSE terendah ($1190.549$) dan koefisien korelasi tertinggi ($0.778$), disusul oleh Random Forest. Sebaliknya, Decision Tree tunggal menghasilkan tingkat galat tertinggi akibat sifatnya yang sangat kaku dan rentan \textit{overfitting}.
\end{itemize}

\clearpage
\section{Tugas Praktikum}
Pada bagian Tugas Praktikum, dilakukan penyelesaian dua studi kasus peramalan harga saham menggunakan \textit{Tesla Stock Data from 2010 to 2020}. Studi kasus ini menguji kemampuan model \textit{machine learning} dalam memprediksi harga penutupan (\textit{Close}) saham 2 hari ke depan ($t+2$) berdasarkan kondisi pergerakan pasar historis 7 hari sebelumnya ($t-7$ hingga $t-1$).

\subsection{Tugas 1: Forecasting Target Close 2 Hari ke Depan (Fitur Close 7 Hari Sebelumnya)}
Studi kasus pertama memodelkan peramalan univariat, di mana nilai target harga \textit{Close} pada horizon 2 hari ke depan diprediksi secara eksklusif menggunakan rekayasa fitur dari data deret waktu harga \textit{Close} selama 7 hari sebelumnya.

\subsubsection{Import Library dan Pemuatan Dataset Tesla Stock}
Langkah awal meliputi inisialisasi lingkungan komputasi dan pemuatan berkas data historis pergerakan saham Tesla.

\begin{figure}[H]
    \centering
    \includegraphics[width=0.35\textwidth]{T1-1,2a.png}
    \caption{Source Code Import Library dan Load Dataset Tesla Stock}
    \label{fig:t1_1}
\end{figure}
Penjelasan Gambar :
\begin{itemize}
    \item \textbf{Input:} Mengimpor pustaka pengolahan matriks dan analisis deret waktu, kemudian membaca berkas \texttt{TSLA.csv} menggunakan fungsi \texttt{pd.read\_csv()}. Inspeksi struktur kolom dilakukan menggunakan \texttt{df.head()}.
    \item \textbf{Output:} Menampilkan cuplikan 5 baris awal dataset saham Tesla yang terdiri dari kolom \texttt{Date}, \texttt{Open}, \texttt{High}, \texttt{Low}, \texttt{Close}, \texttt{Adj Close}, dan \texttt{Volume} sejak tahun 2010.
\end{itemize}

\subsubsection{Penyesuaian Format Data Target (Close)}
Data harga penutupan diisolasi dan dikonversi ke dalam format larik numerik kontinu.

\begin{figure}[H]
    \centering
    \includegraphics[width=0.35\textwidth]{T1-2b.png}
    \caption{Source Code Isolasi Fitur Deret Waktu Close}
    \label{fig:t1_2}
\end{figure}
Penjelasan Gambar :
\begin{itemize}
    \item \textbf{Input:} Mengekstrak kolom \texttt{Close} dari DataFrame utama, mengonversinya ke dalam tipe data desimal \textit{float}, dan membentuk larik 2 dimensi untuk merepresentasikan runtun waktu harga penutupan.
    \item \textbf{Output:} Larik data numerik harga \textit{Close} yang siap diproses oleh algoritma pembentukan jendela waktu geser.
\end{itemize}

\subsubsection{Implementasi Sliding Window dan Fitur Statistik}
Fungsi modular \textit{sliding window} diimplementasikan dengan parameter jendela masukan 7 hari dan target luaran 2 hari ke depan, diperkaya dengan ekstraksi 8 parameter statistik deskriptif.

\begin{figure}[H]
    \centering
    \includegraphics[width=0.40\textwidth]{T1-3,4.png}
    \caption{Source Code Definisi Sliding Window dan Ekstraksi Fitur Statistik}
    \label{fig:t1_3}
\end{figure}
Penjelasan Gambar :
\begin{itemize}
    \item \textbf{Input:} Menetapkan parameter horizon $n\_steps\_in = 7$ dan $n\_steps\_out = 2$ pada fungsi \texttt{split\_sequences}. Selain itu, fungsi \texttt{stats\_features} disiapkan untuk menghitung statistik dinamis (min, max, selisih rentang, standar deviasi, rata-rata, median, kurtosis, dan skewness) pada setiap jendela transaksi mingguan saham.
    \item \textbf{Output:} Pembentukan arsitektur ekstraksi data deret waktu yang siap memetakan fluktuasi harga 7 hari sebelumnya menuju target harga 2 hari mendatang.
\end{itemize}

\subsubsection{Prapemrosesan Data dan Pembagian Dataset}
Rangkaian prapemrosesan dijalankan untuk mentransformasikan dimensi data, memisahkan data latih dan uji secara kronologis, serta menyematkan fitur statistik.
\begin{figure}[H]
    \centering
    \includegraphics[width=0.45\textwidth]{T1-5a.png}
    \caption{Source Code Prapemrosesan Data dan Ekstraksi Fitur Statistik}
    \label{fig:t1_5a}
    
    \vspace{0.1cm}
    
    \includegraphics[width=0.45\textwidth]{T1-5b.png}
    \caption{Output Verifikasi Matriks Fitur dan Target Latih-Uji}
    \label{fig:t1_5b}
\end{figure}
\begin{figure}[H]
    \centering
    \includegraphics[width=0.65\textwidth]{T1-5c.png}
    \caption{Plot Visualisasi Tren Pergerakan Historis Harga Saham Tesla}
    \label{fig:t1_5c}
\end{figure}
Penjelasan Gambar :
\begin{itemize}
    \item \textbf{Input:} Pada Gambar \ref{fig:t1_5a}, data deret waktu harga penutupan (\texttt{in\_seq}) dipecah menjadi pasangan jendela masukan dan keluaran menggunakan fungsi \texttt{split\_sequences(in\_seq, 7, 2)}, kemudian diratakan dimensinya (\texttt{X.reshape}) dan dibagi menjadi data latih (80\%) serta data uji (20\%) melalui \texttt{train\_test\_split(..., test\_size=0.2, shuffle=False)} guna mempertahankan urutan temporal data. Selanjutnya, 8 fitur statistik deskriptif disematkan ke dalam \texttt{X\_train} dan \texttt{X\_test} menggunakan fungsi \texttt{stats\_features}. Pada Gambar \ref{fig:t1_5c}, dieksekusi perintah visualisasi grafik runtun waktu harga penutupan saham Tesla menggunakan pustaka \texttt{matplotlib.pyplot}.
    \item \textbf{Output:} Gambar \ref{fig:t1_5b} memvalidasi bentuk matriks dan nilai sampel: data uji memiliki 483 baris sampel dengan 15 fitur (7 fitur lag harga \textit{Close} ditambah 8 fitur ringkasan statistik) serta target harga 2 hari ke depan. Gambar \ref{fig:t1_5c} menampilkan grafik tren pergerakan harga saham Tesla dari rentang tahun 2010 hingga 2020 yang memperlihatkan dinamika harga relatif stabil pada periode awal dan mengalami tren kenaikan eksponensial yang tajam (\textit{bullish}) pada periode 2019--2020.
\end{itemize}

\subsubsection{Pelatihan Model Machine Learning}
Empat model regresi dilatih menggunakan data latih pergerakan harga saham Tesla.

\begin{figure}[H]
    \centering
    \includegraphics[width=0.65\textwidth]{T1-6.png}
    \caption{Source Code Eksekusi Pelatihan Empat Model Machine Learning}
    \label{fig:t1_6}
\end{figure}
Penjelasan Gambar :
\begin{itemize}
    \item \textbf{Input:} Memanggil fungsi pelatihan \texttt{mlp()}, \texttt{knn()}, \texttt{dt()}, dan \texttt{rf()} dengan melewatkan matriks \texttt{X\_train}, \texttt{X\_test}, \texttt{y\_train}, dan \texttt{y\_test} berbasis fitur harga \textit{Close}.
    \item \textbf{Output:} Keempat model berhasil mempelajari pola fluktuasi harga saham dan menghasilkan parameter internal untuk melakukan inferensi.
\end{itemize}

\subsubsection{Pengujian dan Prediksi}
Tahap inferensi dijalankan untuk menghasilkan taksiran harga saham pada periode uji.

\begin{figure}[H]
    \centering
    \includegraphics[width=0.65\textwidth]{T1-7.png}
    \caption{Source Code Inferensi dan Pembulatan Prediksi}
    \label{fig:t1_7}
\end{figure}
Penjelasan Gambar :
\begin{itemize}
    \item \textbf{Input:} Eksekusi komputasi fungsi \texttt{predict()} dari model terhadap matriks uji \texttt{X\_test} serta pembulatan nilai hasil peramalan.
    \item \textbf{Output:} Vektor nilai prediksi harga penutupan saham Tesla 2 hari ke depan untuk tiap titik waktu pengujian.
\end{itemize}

\clearpage
\subsubsection{Visualisasi Hasil Prediksi Model}
Perbandingan kurva hasil estimasi harga saham terhadap pergerakan harga saham aktual dievaluasi secara grafis.

\begin{figure}[H]
    \centering
    \includegraphics[width=0.50\textwidth]{T1-8a.png}
    \caption{Plot Visualisasi Hasil Prediksi Model MLP dan KNN pada Tugas 1}
    \label{fig:t1_8a}
    
    \vspace{0.1cm}
    
    \includegraphics[width=0.50\textwidth]{T1-8b.png}
    \caption{Plot Visualisasi Hasil Prediksi Model DT dan RF pada Tugas 1}
    \label{fig:t1_8b}
\end{figure}
Penjelasan Gambar :
\begin{itemize}
    \item \textbf{Input:} Perintah plotting kurva deret waktu harga penutupan riil bersanding dengan estimasi model MLP dan KNN (Gambar \ref{fig:t1_8a}), serta model Decision Tree dan Random Forest (Gambar \ref{fig:t1_8b}).
    \item \textbf{Output:} Gambar \ref{fig:t1_8a} memperlihatkan bahwa kurva MLP dan KNN mampu mengikuti pergerakan umum harga saham, meskipun pada lonjakan harga ekstrem terdapat jarak galat yang nyata. Gambar \ref{fig:t1_8b} menunjukkan bahwa Random Forest menghasilkan kurva prediksi yang paling mulus dan konsisten mendekati data aktual dibandingkan Decision Tree yang memiliki deviasi tajam pada titik balik pasar.
\end{itemize}

\subsubsection{Evaluasi Performa Model Tugas 1}
Hasil evaluasi kuantitatif dari keempat model pada penugasan peramalan univariat diukur menggunakan RMSE dan korelasi Pearson.

\begin{figure}[H]
    \centering
    \includegraphics[width=0.6\textwidth]{T1-9.png}
    \caption{Output Evaluasi Metrik Performa Model Tugas 1}
    \label{fig:t1_9}
\end{figure}
Penjelasan Gambar :
\begin{itemize}
    \item \textbf{Input:} Perintah cetak metrik kesalahan RMSE dan tingkat korelasi Pearson untuk model MLP, KNN, Decision Tree, dan Random Forest pada Tugas 1.
    \item \textbf{Output:} Menampilkan perbandingan performa model. Berdasarkan hasil metrik pada Gambar \ref{fig:t1_9}, model berbasis \textit{ensemble learning} \textbf{Random Forest (RF)} dan jaringan saraf tiruan \textbf{MLP} mencatatkan nilai korelasi Pearson yang sangat tinggi ($r > 0.95$), menunjukkan bahwa kedua model sangat handal dalam menangkap tren pergerakan harga saham jangka pendek. Sebaliknya, Decision Tree tunggal menunjukkan kerentanan yang tinggi terhadap volatilitas pasar saham dengan nilai RMSE terbesar.
\end{itemize}

\clearpage
\subsection{Tugas 2: Forecasting Target Close 2 Hari ke Depan (Fitur Close dan Open 7 Hari Sebelumnya)}
Studi kasus kedua menerapkan pendekatan multivariat, di mana variabel input diperkaya dengan memasukkan informasi harga pembukaan (\texttt{Open}) bersama dengan harga penutupan (\texttt{Close}) selama 7 hari sebelumnya ($t-7$ hingga $t-1$) untuk memprediksi harga penutupan 2 hari ke depan ($t+2$).

\subsubsection{Pemuatan Dataset dan Penyesuaian Fitur Multi-Variabel (Open dan Close)}
Data saham diekstrak untuk membentuk matriks runtun waktu multivariat yang menggabungkan dinamika pembukaan dan penutupan bursa.

\begin{figure}[H]
    \centering
    \includegraphics[width=0.55\textwidth]{T2-1.png}
    \caption{Source Code Pemuatan dan Ekstraksi Fitur Multivariat (Open dan Close)}
    \label{fig:t2_1}
\end{figure}
Penjelasan Gambar :
\begin{itemize}
    \item \textbf{Input:} Mengambil kolom \texttt{Open} dan \texttt{Close} dari DataFrame saham Tesla menggunakan sintaks \texttt{df[['Open', 'Close']].astype(float).values} untuk menghasilkan matriks masukan multivariat dua dimensi.
    \item \textbf{Output:} Matriks numerik berdimensi $N \times 2$ yang memuat riwayat harga pembukaan dan penutupan harian secara simultan.
\end{itemize}

\subsubsection{Implementasi Sliding Window Multi-Variabel dan Fitur Statistik}
Fungsi transformasi deret waktu diadaptasi untuk menangani jendela masukan multivariat secara konsisten.

\begin{figure}[H]
    \centering
    \includegraphics[width=0.55\textwidth]{T2-2,3.png}
    \caption{Source Code Definisi Sliding Window Multivariat dan Fitur Statistik}
    \label{fig:t2_2}
\end{figure}
Penjelasan Gambar :
\begin{itemize}
    \item \textbf{Input:} Menjalankan fungsi \texttt{split\_sequences} pada matriks multivariat dengan parameter jendela 7 langkah input dan 2 langkah target output penutupan, disertai perhitungan ringkasan statistik deskriptif pada tiap blok waktu.
    \item \textbf{Output:} Fungsi transformasi multivariat yang siap memetakan interaksi harga pembukaan dan penutupan mingguan ke target harga penutupan masa depan.
\end{itemize}

\subsubsection{Prapemrosesan Data dan Pelatihan Model Machine Learning}
Pipeline prapemrosesan dijalankan untuk membentuk dataset latih dan uji multivariat, disusul pelatihan keempat model regresi.

\begin{figure}[H]
    \centering
    \includegraphics[width=0.50\textwidth]{T2-4,5.png}
    \caption{Source Code Prapemrosesan Multivariat dan Pelatihan Model}
    \label{fig:t2_4}
\end{figure}
Penjelasan Gambar :
\begin{itemize}
    \item \textbf{Input:} Membagi matriks multivariat menjadi data latih dan uji secara kronologis (\texttt{shuffle=False}), menyematkan fitur statistik, dan melatih model \texttt{mlp()}, \texttt{knn()}, \texttt{dt()}, serta \texttt{rf()}.
    \item \textbf{Output:} Matriks masukan multivariat yang telah diperkaya statistik berhasil dipelajari oleh keempat algoritma \textit{machine learning}.
\end{itemize}

\subsubsection{Pengujian dan Prediksi}
Eksekusi pengujian dilakukan untuk menghasilkan nilai estimasi harga penutupan saham Tesla berbasis model multivariat.

\begin{figure}[H]
    \centering
    \includegraphics[width=0.50\textwidth]{T2-6.png}
    \caption{Source Code Inferensi Model Multivariat}
    \label{fig:t2_6}
\end{figure}
Penjelasan Gambar :
\begin{itemize}
    \item \textbf{Input:} Pemanggilan fungsi prediksi dari masing-masing model terhadap matriks uji multivariat \texttt{X\_test}.
    \item \textbf{Output:} Vektor numerik hasil peramalan harga saham penutupan 2 hari ke depan pada pengujian multivariat.
\end{itemize}

\subsubsection{Visualisasi Hasil Prediksi Model}
Visualisasi komparatif kurva peramalan multivariat dipetakan terhadap pergerakan data aktual.

\begin{figure}[H]
    \centering
    \includegraphics[width=0.50\textwidth]{T2-7a.png}
    \caption{Plot Visualisasi Hasil Prediksi Model MLP dan KNN pada Tugas 2}
    \label{fig:t2_7a}
    
    \vspace{0.3cm}
    
    \includegraphics[width=0.50\textwidth]{T2-7b.png}
    \caption{Plot Visualisasi Hasil Prediksi Model DT dan RF pada Tugas 2}
    \label{fig:t2_7b}
\end{figure}
Penjelasan Gambar :
\begin{itemize}
    \item \textbf{Input:} Script plotting \texttt{matplotlib} yang memetakan nilai aktual \texttt{y\_test} bersanding dengan hasil peramalan model multivariat MLP dan KNN (Gambar \ref{fig:t2_7a}) serta model DT dan RF (Gambar \ref{fig:t2_7b}).
    \item \textbf{Output:} Gambar \ref{fig:t2_7a} dan Gambar \ref{fig:t2_7b} memperlihatkan bahwa penyertaan fitur harga \texttt{Open} secara signifikan membantu model dalam mendeteksi perubahan sentimen harian bursa. Model Random Forest dan MLP memperlihatkan adaptasi yang sangat baik dalam merekonstruksi lonjakan harga (\textit{price surge}) pada akhir periode data uji.
\end{itemize}

\subsubsection{Evaluasi Performa Model Tugas 2}
Evaluasi metrik kuantitatif dianalisis untuk memverifikasi dampak penambahan fitur harga pembukaan (\textit{Open}).

\begin{figure}[H]
    \centering
    \includegraphics[width=0.55\textwidth]{T2-8.png}
    \caption{Output Evaluasi Metrik Performa Model Tugas 2}
    \label{fig:t2_8}
\end{figure}
Penjelasan Gambar :
\begin{itemize}
    \item \textbf{Input:} Perintah pencetakan metrik evaluasi RMSE dan korelasi Pearson untuk keempat model pada Tugas 2.
    \item \textbf{Output:} Menampilkan log performa numerik model multivariat. Berdasarkan Gambar \ref{fig:t2_8}, integrasi fitur harga \texttt{Open} bersama harga \texttt{Close} menghasilkan peningkatan akurasi peramalan dan stabilitas metrik pada model ensemble \textbf{Random Forest} dan \textbf{MLP}, yang dibuktikan dengan penurunan nilai galat RMSE dan peningkatan konsistensi korelasi linear mendekati nilai $1.0$.
\end{itemize}

\subsection{Analisis Komparatif: Dampak Penambahan Fitur Open (Tugas 1 vs Tugas 2)}
Perbandingan antara eksperimen Tugas 1 (univariat \textit{Close}) dan Tugas 2 (multivariat \textit{Open} dan \textit{Close}) memberikan wawasan penting dalam analisis peramalan data finansial:
\begin{enumerate}
    \item \textbf{Pengayaan Informasi Pasar:} Pada Tugas 1, model hanya melihat harga akhir transaksi harian tanpa mengetahui dinamika pembukaan bursa. Pada Tugas 2, penambahan fitur harga pembukaan (\texttt{Open}) memberikan informasi mengenai \textit{gap} harga (selisih antara harga penutupan hari sebelumnya dan pembukaan hari berikutnya), yang sering merefleksikan sentimen pasar di luar jam perdagangan aktif.
    \item \textbf{Stabilitas Prediksi Algoritma Ensemble:} Algoritma \textit{Random Forest} menunjukkan peningkatan ketahanan yang paling nyata saat fitur bertambah. Kemampuannya melakukan pemilihan subset fitur secara acak (\textit{feature bagging}) memungkinkan model memanfaatkan interaksi antara harga \textit{Open} dan \textit{Close} tanpa mengalami \textit{overfitting}.
    \item \textbf{Kelemahan Model Kaku:} Baik pada Tugas 1 maupun Tugas 2, algoritma \textit{Decision Tree} tunggal tetap mencatatkan performa terlemah. Hal ini membuktikan bahwa pembagian ruang fitur berbasis ambang batas kaku (\textit{orthogonal splits}) tidak mampu menangkap tren deret waktu yang bersifat kontinu dan dinamis secara efektif.
\end{enumerate}

\chapter{Kesimpulan}
\label{chap:kesimpulan}
Berdasarkan praktikum mengenai \textit{Forecasting} (Peramalan Deret Waktu) berbasis \textit{machine learning} yang telah dilaksanakan—meliputi pemodelan jumlah peminjaman sepeda pada \textit{Bike Sharing Dataset} serta peramalan univariat dan multivariat harga saham pada \textit{Tesla Stock Data}—diperoleh beberapa simpulan penting sebagai berikut:

\begin{enumerate}
    \item \textbf{Efektivitas Rekayasa Fitur Sliding Window dan Agregasi Statistik:} Metode \textit{sliding window} terbukti sukses mentransformasikan struktur deret waktu sekuensial menjadi format pasangan data berlabel yang dapat dipelajari oleh algoritma \textit{supervised learning}. Selain itu, penambahan 8 fitur statistik deskriptif lokal (\textit{mean, median, variance, std dev, min, max, kurtosis,} dan \textit{skewness}) mampu memperkaya representasi volatilitas data dan meningkatkan akurasi estimasi secara signifikan tanpa harus memperpanjang jendela lag secara berlebihan.
    
    \item \textbf{Keunggulan Algoritma Ensemble dan Jaringan Saraf Tiruan:} Model berbasis ansambel (\textit{Random Forest}) dan jaringan saraf tiruan (\textit{Multilayer Perceptron}) secara konsisten mencatatkan performa tertinggi dalam menangkap pola tren non-linear. Mekanisme agregasi pohon pada \textit{Random Forest} memberikan stabilitas luar biasa terhadap derau (\textit{noise}) pasar, sedangkan fleksibilitas fungsi aktivasi non-linear pada MLP menghasilkan kurva proyeksi yang sangat adaptif terhadap fluktuasi data riil.
    
    \item \textbf{Kelemahan Inheren Model Pohon Keputusan Tunggal dan KNN:} Algoritma \textit{Decision Tree} tunggal terbukti menghasilkan nilai kesalahan galat (RMSE) tertinggi dan korelasi terendah pada seluruh eksperimen akibat sifat partisi ortogonalnya yang kaku dan rentan terhadap \textit{overfitting}. Sementara itu, algoritma KNN menunjukkan keterbatasan dalam melakukan ekstrapolasi tren ekstrem di luar rentang sampel data latih (\textit{extrapolation limitation}).
    
    \item \textbf{Dampak Positif Fitur Multivariat dalam Peramalan Saham:} Pengujian pada dataset saham Tesla membuktikan bahwa penggabungan fitur harga pembukaan (\texttt{Open}) bersama harga penutupan (\texttt{Close}) pada Tugas 2 menghasilkan estimasi yang lebih presisi dibandingkan peramalan univariat pada Tugas 1. Informasi gap harga antara pembukaan dan penutupan harian memberikan sinyal krusial terkait dinamika sentimen pasar yang memperkuat kemampuan inferensi model.
    
    \item \textbf{Pentingnya Kombinasi Metrik RMSE dan Koefisien Korelasi Pearson:} Evaluasi model peramalan deret waktu tidak cukup hanya diukur dari selisih magnitudo galat kuadrat (RMSE), melainkan wajib dipadukan dengan Koefisien Korelasi Pearson ($r$). Kombinasi ini memberikan gambaran komprehensif, di mana RMSE mengukur presisi skala numerik sedangkan korelasi Pearson memverifikasi keakuratan arah tren pergerakan pasar.
\end{enumerate}

\chapter{Daftar Pustaka}
\begin{enumerate}
    \item \label{modulPraktikum} Dr. Eng. Ir. Ganjar Alfian, S.T., M. Eng., dan Tim Asisten. \textit{Modul Praktikum Penambangan Data: Pertemuan V Forecasting}. Yogyakarta: Departemen Teknik Elektro dan Informatika, Sekolah Vokasi, Universitas Gadjah Mada, 2025.
    
    \item \label{retnoningsih} E. Retnoningsih and R. Pramudita, "Mengenal Machine Learning Dengan Teknik Supervised dan Unsupervised Learning Menggunakan Python," \textit{Bina Insani ICT Journal}, vol. 7, no. 2, pp. 156--165, 2020. [Online]. Tersedia: \url{https://doi.org/10.51211/biict.v7i2.1422}.
    
    \item \label{santoso} H. Abijono, P. Santoso, and N. L. Anggreini, "Algoritma Supervised Learning dan Unsupervised Learning dalam Pengolahan Data," \textit{G-Tech: Jurnal Teknologi Terapan}, vol. 4, no. 2, pp. 315--318, 2021. [Online]. Tersedia: \url{https://doi.org/10.33379/gtech.v4i2.635}.
    
    \item \label{vimala} J. Vimala and A. Nugroho, "Forecasting Penjualan Obat Menggunakan Metode Single, Double, dan Triple Exponential Smoothing (Studi Kasus: Apotek Mandiri Medika)," \textit{IT-Explore: Jurnal Penerapan Teknologi Informasi dan Komunikasi}, vol. 1, no. 2, pp. 90--99, 2022. [Online]. Tersedia: \url{https://ejournal.uksw.edu/it-explore/article/view/5888}.
    
    \item \label{maysofa} R. Maysofa, I. Syaliman, and D. Sapriadi, "Implementasi Forecasting pada Penjualan Inaura Hair Care dengan Metode Single Exponential Smoothing," \textit{Jurnal Testing dan Implementasi Sistem Informasi}, vol. 1, no. 2, pp. 82--91, 2023. [Online]. Tersedia: \url{https://journal.al-matani.com/index.php/jtisi/article/view/423}.
    
    \item \label{chai} T. Chai and R. R. Draxler, "Root mean square error (RMSE) or mean absolute error (MAE)? Arguments against avoiding RMSE in the literature," \textit{Geoscientific Model Development}, vol. 7, no. 3, pp. 1247--1250, 2014. Tersedia: \url{https://doi.org/10.5194/gmd-7-1247-2014}.
    
    \item \label{benesty} J. Benesty, J. Chen, Y. Huang, and I. Cohen, "Pearson Correlation Coefficient," in \textit{Noise Reduction in Speech Processing}, Berlin, Heidelberg: Springer, 2009, pp. 1--4. [Online]. Tersedia: \url{https://doi.org/10.1007/978-3-642-00296-0_5}.
    
    \item \label{kaggleBikesharing} L. N. Pathi, \textit{Bike Sharing Dataset}. Kaggle, 2018. [Online]. Tersedia: \url{https://www.kaggle.com/datasets/lakshmi25npathi/bike-sharing-dataset}.
    
    \item \label{kaggleTesla} T. Boz, \textit{Tesla Stock Data from 2010 to 2020}. Kaggle, 2020. [Online]. Tersedia: \url{https://www.kaggle.com/datasets/timoboz/tesla-stockdata-from-2010-to-2020/data}.
    
    \item \label{colabLaporan} H. C. A. Paliling, \textit{S5-PPD-542159 (Source Code Laporan Praktikum Pertemuan 5 - Forecasting)}. Google Colaboratory, 2026. [Online]. Tersedia: \url{https://colab.research.google.com/drive/1example_link_colab_ppd5}.
\end{enumerate}

\end{document}